
\section[Ser cohomológicamente propio y aciclicidad local]{Propiedad de ser cohomológicamente propio y aciclicidad local genérica}
\label{XIII.3}
\marginpar{395}

\begin{theoreme}
\label{XIII.3.1}
Sean $S$ un esquema irreducible con punto genérico~$s$, $X$ e~$Y$ dos
$S$-esquemas de presentación finita y $f\colon X\to Y$ un
$S$-morfismo. Para todo $S$-esquema~$S'$, se denotan por $Y'$, $X'$, etc.
las imágenes inversas de~$Y$, $X$, etc. por el morfismo $S'\to S$.
Se tienen las propiedades siguientes:

\begin{enumerate}
\item[1)]
{\upshape a)} Puede encontrarse un abierto no vacío~$S'$ de~$S$ tal que,
para todo haz de conjuntos constante finito~$F'$ sobre~$X'$, $f'_*F'$ sea
constructible y~$(F',f')$ sea cohomológicamente propio relativamente
a~$S'$ en dimensión $\leqslant 0$.

{\upshape b)} Sea $F$ un haz de conjuntos constructible sobre~$X$.
Entonces puede encontrarse un abierto no vacío~$S'$ de~$S$ (dependiente
de~$F$) tal que $f'_*F'$ sea constructible y~$(F',f')$ sea
cohomológicamente propio relativamente a~$S'$ en dimensión
$\leqslant 0$.

\item[2)] Supongamos que los esquemas de tipo finito de dimensión
$\leqslant\dim X_s$ sobre una clausura algebraica~$\overline{k}$
de~$\kres(s)$ son fuertemente desingularizables
\textup{(SGA~5 I~3.1.5)}. Se tienen entonces, además, las propiedades
siguientes:

{\upshape a)} Puede encontrarse un abierto no vacío~$S'$ de~$S$ tal que,
para todo haz de grupos constante finito~$F'$ sobre~$X'$, de orden primo
con las características residuales de~$S$, si~$\Phi'$ es el stack de los
torsores bajo~$F'$, $f'_*\Phi'$ sea $1$-constructible y~$(F',f')$ sea
cohomológicamente propio relativamente a~$S'$ en dimensión
$\leqslant 1$.

{\upshape b)} Sean $\LL$ el conjunto de los números primos distintos de
las características residuales de~$S$ y~$\Phi$ un stack ind-$\LL$,
$1$-constructible, sobre~$X$ \eqref{XIII.0}, tal que, para todo
esquema~$X_1$ étale sobre~$X$ y para todo par de objetos $x,x_1$
de~$\Phi_{X_1}$, el haz $\SheafHom_{X_1}(x,x_1)$ sea constructible.
Supongamos además que~$S$ es localmente noetheriano. Entonces puede
encontrarse un abierto no vacío~$S'$ de~$S$ tal que $f'_*\Phi'$ sea
$1$-constructible, que,
\ifthenelse{\boolean{orig}}
{}
{para todo esquema~$Y_1$ étale sobre~$Y'$ y}
para todo par de objetos $y,y_1$ de la fibra $(f'_*\Phi')_{Y_1}$,
el haz $\SheafHom_{Y_1}(y,y_1)$ sea constructible,
\marginpar{396}
y que~$(\Phi',f')$ sea cohomológicamente propio relativamente a~$S'$
en dimensión $\leqslant 1$.
\end{enumerate}
\end{theoreme}

\subsubsection*{Demostración}
Puede suponerse que~$S$ es afín; según SGA~4 VIII~1.1, puede suponerse
que~$S$ es íntegro; finalmente, por paso al límite, puede suponerse que
$S$ es el espectro de un álgebra de tipo finito sobre~$\ZZ$; en
particular, $S$ es entonces noetheriano.
Como la cuestión es local sobre~$Y$, puede suponerse que~$Y$ es afín.
Además, para demostrar el teorema, basta hacerlo tras una extensión finita
de la base $S'\to S$, donde~$S'$ es un esquema íntegro y~$S'\to S$ es un
compuesto de morfismos étales y de morfismos finitos, radiciales y
sobreyectivos.
